\par
We interpret $\omega_1$ as the first uncountable ordinal and $\omega_1^2$ as the ordinal product $\omega_1\cdot\omega_1$. Likewise $\omega_1\omega$ denotes $\omega_1\cdot\omega$.

Fix a (simple, undirected) graph $G=(V(G),E(G))$ with $|V(G)|\le\aleph_1$ such that $G$ has no subgraph isomorphic to $K_4$ and no subgraph isomorphic to $K_{\aleph_0,\aleph_0}$.

For ordinals $\alpha,\beta$ and a graph $G$, we take
\[
\alpha\to(\beta,G)^2
\]
to mean: for every graph $H=(\alpha,E_H)$ on vertex set $\alpha$, either
\begin{itemize}
\item there exists an \emph{independent} set $I\subseteq\alpha$ such that $H[I]$ has no edges and the order type of $I$ (with the inherited well-order) is exactly $\beta$, or
\item there exists an injection $f:V(G)\to\alpha$ such that for every edge $\{u,v\}\in E(G)$ we have $\{f(u),f(v)\}\in E_H$ (so $H$ contains a subgraph isomorphic to $G$).
\end{itemize}

\emph{Ambiguity note.} If one instead requires an \emph{induced} copy of $G$ in the second bullet, then even a clique of size $|V(G)|$ would not necessarily contain $G$. The surrounding discussion (``contains no $K_4$'') strongly suggests ``subgraph'' rather than ``induced subgraph'', and we adopt ``subgraph''.

Stress points: the order types $\omega_1^2$ and $\omega_1\omega$; the forbidden subgraphs $K_4$ and $K_{\aleph_0,\aleph_0}$; and whether $G$ is finite vs size $\aleph_1$.
